% !PRINTED 583
\weberpage{583}
%<<CONT
are. If $\varrho$ is a third root of unity, we can set $\varrho = r^3$,
and obtain the resolvents
\weq{24}{(\varrho,\eta) = \eta + \varrho\eta' + \varrho^2\eta'', \quad (\varrho^{-1},\eta) = \eta + \varrho^{-1}\eta' + \varrho^{-2}\eta'',}
for which, by (22), making use of (21), one obtains
\[
(\varrho, r) = 3r, \quad (\varrho^{-1}, r) = 3r^{-1},
\]
hence
\weq{25}{(\varrho, r)^3 = 27\varrho,}
\weq{26}{(\varrho, r)(\varrho^{-1}, r) = 9,}
% [source-original mathematical notation, verified 2026-08-29] Immediately after defining
% the resolvents (varrho,eta) in (24), the print switches to (varrho,r) in the two displayed
% evaluations and in (25)-(26). Direct substitution from (21)-(24) gives (varrho,eta)=3r
% and (varrho^-1,eta)=3r^-1, so eta is the expected second argument; print preserved.
which is entirely analogous to the formulae (5) and (6).

We enter briefly also on the case $e = 4$, which occurs for
the primes $n$ of the form $4f + 1$ ($f$ an integer),
thus for $n = 5, 13, 17, 29, 37, 41 \dots$

Let the four periods of $f$ terms each again be $\eta, \eta_1, \eta_2, \eta_3$.

For $\alpha$ in \S.\,169 (8), we have to set $i$, and thus obtain,
since $-2 \equiv +2\ (\mathrm{mod}\ 4)$,
\weq{27}{\begin{aligned}
\psi_1(i) &= \sum i^{\operatorname{ind}t + 2\operatorname{ind}(t+1)}\\
\psi_2(i) &= \sum i^{\operatorname{ind}(t^2+t)} = a + bi,
\end{aligned}}
wherein $a$ and $b$ are integers. $\psi_1(i)$ has the same form
as $\psi_2(i)$. But we shall presently reduce the one function to the
other.

By the results of \S.\,171 we have
\weq{28}{(-1,\eta) = \eta + \eta_2 - \eta_1 - \eta_3 = \sqrt{n},}
and hence
\weq{29}{2(\eta + \eta_2) = -1 + \sqrt{n}, \quad 2(\eta_1 + \eta_3) = -1 - \sqrt{n},}
where the sign of $\sqrt{n}$ is to be determined according to \S.\,171 and
may, on a suitable assumption about $r$, be taken positive.

We have further, by \S.\,169, (14), (15)
\weq{30}{\begin{aligned}
(i,\eta)(-i,\eta) &= (-1)^f n\\
(i,\eta)^2 &= (-1,\eta)\psi_1(i),
\end{aligned}}
hence by (28)
\weq{31}{(i,\eta)^2 = \psi_1(i)\sqrt{n}, \quad (-i,\eta)^2 = \psi_1(-i)\sqrt{n}}
\weq{32}{(i,\eta)^4 = (-1)^f n\,\psi_1(i)\psi_2(i) = n\,\psi_1(i)^2,}
from which it follows
\weq{33}{\psi_1(i) = (-1)^f \psi_2(i) = (-1)^f(a + bi),}
and from this by (30)
\weq{34}{\psi_1(i)\psi_1(-i) = \psi_2(i)\psi_2(-i) = a^2 + b^2 = n.}
